\usepackage{matplotlibcolors}
% Define an environment for a slide with a background image
\newenvironment{frameb}[2]
{	
	\IfEqCase{#1}{ 	% use the first argument to determine
			% whether to fit to height or to width
			%
	{h}		% fit to height
	{\usebackgroundtemplate{\includegraphics[height=\paperheight]{#2}}}
	{w}		% fit to width 
	{\usebackgroundtemplate{\includegraphics[width=\paperwidth]{#2}}}
			% error message:
	{\PackageError{frameb}{Undefined option to frameb: #1 first argument must either be h to stretch to height, or w to stretch to width}
	}}
	\begin{frame}   
}
{
	\end{frame}  
}

% Define an environment for a slide with a background image drawn by tikz
\newenvironment{framet}[3]
{	
% re-set the background... 
\setbeamertemplate{background}{
% ... as a TikZ box, with unitz re-scaled to 4:3 units
\begin{tikzpicture}[x=0.25\paperwidth,y=0.33333333\paperheight]
%
% If we want to put an image as well, then you should put:
\IfEqCase{#1}{
{h}    % Fit to height
{
\node[inner sep=0pt, anchor=south west] at (0,0) {\includegraphics[height=\paperheight]{#2}};
}
{w}    % Fit to width
{
\node[inner sep=0pt, anchor=north west] at (0,3) {\includegraphics[width=\paperwidth]{#2}};
}
{n}    % No background image
{}
{\PackageError{framet}{Undefined option to framet: #1 first argument must either be h to stretch to height, w to stretch to width, or n for no background image}
}
}
% Now the rest of the TikZ code
#3
%
% Finish the Tikz background
\end{tikzpicture}
%
}   % done setting the background
%
% Now begin the frame
	\begin{frame}   
}
{
% close off the frame
	\end{frame}  
}
